\section{The Rupture with Alexandria}

\subsection{A layman in the pulpit}

The quarrel that ended Origen's Alexandrian life began with a question of
church order. When ``a considerable war broke out in the city''---commonly
identified with the disturbances under Caracalla about 215/216---Origen,
``thinking that it would be unsafe for him to remain in Egypt,'' withdrew to
Palestine ``and abode in Caesarea,'' where ``the bishops of the church in
that country requested him to preach and expound the Scriptures publicly,
although he had not yet been ordained as presbyter'' (\work{Church History}
6.19.16). Demetrius protested that laymen do not preach before bishops.
Alexander of Jerusalem and Theoctistus of Caesarea answered him in a letter
Eusebius preserves: the claim ``that such a thing was never heard of before
\ldots\ is plainly untrue,'' ``for whenever persons able to instruct the
brethren are found, they are exhorted by the holy bishops to preach to the
people''---and they name precedents from three churches of Asia Minor
(\work{Church History} 6.19.17--18). Demetrius summoned his catechist home
``through members and deacons of the church,'' and Origen went (\work{Church
History} 6.19.19). The exchange displays the fault line a decade before the
break: two Palestinian bishops who treated Origen as a charism to be used,
and an Alexandrian bishop who treated him as a subordinate to be governed.

These same years brought imperial attention. While the court of Alexander
Severus was at Antioch, the emperor's mother Julia Mamaea, ``a most pious
woman, if there ever was one,'' ``sent for him with a military escort'' and
heard him expound ``the glory of the Lord and of the excellence of the
divine teaching'' before he ``hastened back to his accustomed work''
(\work{Church History} 6.21.3--4).

\subsection{Ordination, condemnation, departure}

About 231 Origen was sent to Greece ``on account of a pressing necessity in
connection with ecclesiastical affairs, and went through Palestine, and was
ordained as presbyter in Caesarea by the bishops of that country''
(\work{Church History} 6.23.4). For Demetrius this was the breaking point: a
foreign ordination of his own catechist, without his leave, and---as his
letters now insisted---of a man allegedly self-mutilated and therefore, on
one reading of church custom, irregular for orders. Eusebius, having told
the mutilation story precisely to expose the bishop's motives, declines to
narrate the proceedings, saying only that ``the matters that were agitated
concerning him on this account, and the decisions on these matters by those
who presided over the churches'' would demand a separate treatise
(\work{Church History} 6.23.4)---the treatise being, in effect, the lost
second book of the \work{Defense of Origen}. Photius, who read that book in
the ninth century, supplies the canonical detail: ``A synod of bishops and
some presbyters was summoned to condemn Origen,'' which ``decided that he
must not remain in Alexandria or teach there, but that he should be allowed
to retain his priesthood. But Demetrius and some Egyptian bishops\ldots\
also deprived him of his sacred office'' (\work{Bibliotheca}, codex 118)
---two stages: banishment from Alexandria, then deposition from the
presbyterate. Jerome, no friend of Origen's doctrine by the time he wrote,
adds the two facts that frame every later assessment: the sentence did not
travel everywhere, and it was not about heresy. ``He stands condemned by his
bishop, Demetrius, only the bishops of Palestine, Arabia, Phenicia, and
Achaia dissenting. Imperial Rome consents to his condemnation,'' and
this---Jerome insists---``not\ldots\ because of the novelty or heterodoxy
of his doctrines, but because men could not tolerate the incomparable
eloquence and knowledge which, when once he opened his lips, made others
seem dumb'' (Jerome, \work{Letter} 33.5).

\witnesslimit{No acts, charges, or signatures of the Alexandrian proceedings
survive. Everything descends from partisan summaries: Eusebius and Pamphilus
defending Origen, Photius epitomizing them six centuries later, Jerome
moralizing the affair as envy. Even the ordaining bishop's name wobbles in
transmission (Photius's epitome names ``Theotecnus'' where Eusebius names
Theoctistus). That Demetrius held proceedings, that Origen left Alexandria
for good in 231/232, and that Palestine, Arabia, Phoenicia, and Achaia never
accepted the sentence are secure; the legal theory, the votes, and the exact
role of the mutilation charge are not. This study asserts no more.}

``It was in the tenth year of the above-mentioned reign''---of Alexander
Severus, 231/232---``that Origen removed from Alexandria to Caesarea,
leaving the charge of the catechetical school in that city to Heraclas. Not
long afterward Demetrius, bishop of the church of Alexandria, died, having
held the office for forty-three full years, and Heraclas succeeded him''
(\work{Church History} 6.26). The sentence outlived its author: Alexandria
under Heraclas and his successor Dionysius---both Origen's former
pupils---never recalled the master. Origen was a presbyter of Caesarea for
the rest of his life.

\subsection{Caesarea: the second school}

In Caesarea, under the protection of Theoctistus, Origen rebuilt everything:
a school, a library---the collection Pamphilus would later curate and
Eusebius inherit---and a preaching ministry, since the ex-catechist was now
a presbyter expounding Scripture in the liturgy. The most vivid surviving
picture of him anywhere comes from this school. Gregory---later called
Thaumaturgus, ``the wonder-worker,'' apostle of Pontus---came to Caesarea
about 233 intending to study Roman law, met Origen, and stayed five years;
his farewell oration describes a pedagogy that began with logic, natural
science, geometry, and astronomy, passed through the ethics of all the
philosophical schools, and ended in Scripture. ``Like some spark lighting
upon our inmost soul, love was kindled and burst into flame within us---a
love at once to the Holy Word\ldots\ and to this man, His friend and
advocate'' (\work{Oration and Panegyric Addressed to Origen} 6). ``No
subject of speech'' was ``forbidden''; ``there was no matter of knowledge
hidden or inaccessible to us, but we had it in our power to learn every kind
of discourse, both foreign and Greek, both spiritual and political, both
divine and human'' (\work{Panegyric} 15). Leaving, Gregory said, felt like
Adam's exile: ``I am like a second Adam and have begun to talk, outside of
paradise'' (\work{Panegyric} 16). Eusebius confirms the brothers Theodore
(Gregory) and Athenodorus as the school's most celebrated alumni, both
afterwards bishops in Pontus (\work{Church History} 6.30). Panegyric is not
a syllabus, and gratitude is not audit; but as evidence of what it felt like
to be taught by Origen, nothing else in ancient Christian literature
compares.

The Caesarean decades also made Origen the East's consulting theologian.
When Beryllus, bishop of Bostra in Arabia, taught that the Saviour had no
distinct pre-existence, a synod called Origen in; he ``went down at first
for a conference with him to ascertain his real opinion,'' then, ``having
persuaded him by argument, and convinced him by demonstration, he brought
him back to the true doctrine'' (\work{Church History} 6.33). When another
Arabian party held that the soul dies with the body until the resurrection,
Origen was summoned again, and again the error was renounced (\work{Church
History} 6.37). The pattern is worth noticing against his later reputation:
in his lifetime, Origen was the man synods called to \emph{end} heresies,
by persuasion rather than anathema.
